% GENERATED FILE. Edit canonical lessons, then run npm run generate:books.
% canonical-source-hash: fnv1a64:a341a56cfe262000
% canonical-lessons: UR-C142-utarna, UR-C142-bharna, UR-C142-chhupna, UR-C142-bachna, UR-C142-barhna, UR-R142-first-pass-describing-words, UR-R142-second-pass-verbs, UR-W10-a1-timed-production

\chapter{To get off, To fill, To hide, To be saved, To grow, A1 timed writing}
\label{ch:ur-to-get-off-to-fill-to-hide-to-be-saved-to-grow}
\hlchaptermodality{142}

\begin{chapteropening}
\textbf{By the end of this chapter:} \emph{I can say to get off, to fill, to hide, to be saved and to grow, then complete the A1 writing paper within 20 minutes.}

Complete the last lesson of chapter 142: I can fill a six-field form and write a 30–40 word message within the A1 paper's 20-minute limit.
\end{chapteropening}

\section[\texorpdfstring{\ur{اترنا} (utarnā)}{اترنا (utarnā)}]{\ur{اترنا} (utarnā) — to get off, to go down}

\label{lesson:UR-C142-utarna}



\begin{quote}
Before the new one: say the Urdu for to come out, then the Urdu for to climb.
\end{quote}

\subsection*{You'll want to know: \ur{اترنا}}
\textbf{\ur{اترنا}} — \emph{utarnā} — \textquotedblleft{}to get off, to go down\textquotedblright{}.

\begin{grammarlens}[title={What you've built}]
One more word for this chapter.
\end{grammarlens}

\subsection*{Your turn}
\begin{itemize}
\raggedright
  \item \emph{Say it:} \emph{utarnā}
  \item \emph{Say it:} \emph{utarnā}, once more
  \item \emph{Say it:} say \emph{utarnā}
  \item \emph{Recall:} say the Urdu for to come out, then the Urdu for to climb, then say \emph{utarnā} again
\end{itemize}

\begin{tcolorbox}[breakable,colback=teal!4,colframe=teal!35!black,title={Before you move on}]
What does \ur{اترنا} mean? (\textquotedblleft{}To get off, to go down\textquotedblright{}.) Say it once more.
\end{tcolorbox}
\section[\texorpdfstring{\ur{بھرنا} (bharnā)}{بھرنا (bharnā)}]{\ur{بھرنا} (bharnā) — to fill}

\label{lesson:UR-C142-bharna}



\begin{quote}
Before the new one: say the Urdu for to climb, then the Urdu for to get off.
\end{quote}

\subsection*{You'll want to know: \ur{بھرنا}}
\textbf{\ur{بھرنا}} — \emph{bharnā} — \textquotedblleft{}to fill\textquotedblright{}.

\begin{grammarlens}[title={What you've built}]
One more word for this chapter.
\end{grammarlens}

\subsection*{Your turn}
\begin{itemize}
\raggedright
  \item \emph{Say it:} \emph{bharnā}
  \item \emph{Say it:} \emph{bharnā}, once more
  \item \emph{Say it:} say \emph{bharnā}
  \item \emph{Recall:} say the Urdu for to climb, then the Urdu for to get off, then say \emph{bharnā} again
\end{itemize}

\begin{tcolorbox}[breakable,colback=teal!4,colframe=teal!35!black,title={Before you move on}]
What does \ur{بھرنا} mean? (\textquotedblleft{}To fill\textquotedblright{}.) Say it once more.
\end{tcolorbox}
\section[\texorpdfstring{\ur{چھپنا} (chhupnā)}{چھپنا (chhupnā)}]{\ur{چھپنا} (chhupnā) — to hide}

\label{lesson:UR-C142-chhupna}



\begin{quote}
Before the new one: say the Urdu for to get off, then the Urdu for to fill.
\end{quote}

\subsection*{You'll want to know: \ur{چھپنا}}
\textbf{\ur{چھپنا}} — \emph{chhupnā} — \textquotedblleft{}to hide\textquotedblright{}.

\begin{grammarlens}[title={What you've built}]
One more word for this chapter.
\end{grammarlens}

\subsection*{Your turn}
\begin{itemize}
\raggedright
  \item \emph{Say it:} \emph{chhupnā}
  \item \emph{Say it:} \emph{chhupnā}, once more
  \item \emph{Say it:} say \emph{chhupnā}
  \item \emph{Recall:} say the Urdu for to get off, then the Urdu for to fill, then say \emph{chhupnā} again
\end{itemize}

\begin{tcolorbox}[breakable,colback=teal!4,colframe=teal!35!black,title={Before you move on}]
What does \ur{چھپنا} mean? (\textquotedblleft{}To hide\textquotedblright{}.) Say it once more.
\end{tcolorbox}
\section[\texorpdfstring{\ur{بچنا} (bachnā)}{بچنا (bachnā)}]{\ur{بچنا} (bachnā) — to be saved, to avoid}

\label{lesson:UR-C142-bachna}



\begin{quote}
Before the new one: say the Urdu for to fill, then the Urdu for to hide.
\end{quote}

\subsection*{You'll want to know: \ur{بچنا}}
\textbf{\ur{بچنا}} — \emph{bachnā} — \textquotedblleft{}to be saved, to avoid\textquotedblright{}.

\begin{grammarlens}[title={What you've built}]
One more word for this chapter.
\end{grammarlens}

\subsection*{Your turn}
\begin{itemize}
\raggedright
  \item \emph{Say it:} \emph{bachnā}
  \item \emph{Say it:} \emph{bachnā}, once more
  \item \emph{Say it:} say \emph{bachnā}
  \item \emph{Recall:} say the Urdu for to fill, then the Urdu for to hide, then say \emph{bachnā} again
\end{itemize}

\begin{tcolorbox}[breakable,colback=teal!4,colframe=teal!35!black,title={Before you move on}]
What does \ur{بچنا} mean? (\textquotedblleft{}To be saved, to avoid\textquotedblright{}.) Say it once more.
\end{tcolorbox}
\section[\texorpdfstring{\ur{بڑھنا} (baṛhnā)}{بڑھنا (baṛhnā)}]{\ur{بڑھنا} (baṛhnā) — to grow, to increase}

\label{lesson:UR-C142-barhna}



\begin{quote}
Before the new one: say the Urdu for to hide, then the Urdu for to be saved.
\end{quote}

\subsection*{You'll want to know: \ur{بڑھنا}}
\textbf{\ur{بڑھنا}} — \emph{baṛhnā} — \textquotedblleft{}to grow, to increase\textquotedblright{}.

\begin{grammarlens}[title={What you've built}]
One more word for this chapter.
\end{grammarlens}

\subsection*{Your turn}
\begin{itemize}
\raggedright
  \item \emph{Say it:} \emph{baṛhnā}
  \item \emph{Say it:} \emph{baṛhnā}, once more
  \item \emph{Say it:} say \emph{baṛhnā}
  \item \emph{Recall:} say the Urdu for to hide, then the Urdu for to be saved, then say \emph{baṛhnā} again
\end{itemize}

\begin{tcolorbox}[breakable,colback=teal!4,colframe=teal!35!black,title={Before you move on}]
What does \ur{بڑھنا} mean? (\textquotedblleft{}To grow, to increase\textquotedblright{}.) Say it once more.
\end{tcolorbox}
\section[(dialogue)]{Describing words — a first pass}

\label{lesson:UR-R142-first-pass-describing-words}



\begin{quote}
Say the words for to be saved and to grow.
\end{quote}

\subsection*{Your turn}
Say these aloud:

\begin{itemize}
\raggedright
  \item young, old, lovely, soft, hard
  \item wet, dry, wide, narrow, deep
  \item famous, brave, clever, tired, hungry
  \item thirsty, fresh, raw, ripe, pink
  \item purple, golden, closed, open, to wake up
\end{itemize}

\begin{tcolorbox}[breakable,colback=teal!4,colframe=teal!35!black,title={Before you move on}]
Young? (\textbf{\ur{جوان}}, \emph{javān}.) Old? (\textbf{\ur{بوڑھا}}, \emph{būṛhā}.) Lovely? (\textbf{\ur{پیارا}}, \emph{pyārā}.) Soft? (\textbf{\ur{نرم}}, \emph{narm}.) Hard? (\textbf{\ur{سخت}}, \emph{sakht}.) Wet? (\textbf{\ur{گیلا}}, \emph{gīlā}.) Dry? (\textbf{\ur{سوکھا}}, \emph{sūkhā}.) Wide? (\textbf{\ur{چوڑا}}, \emph{chauṛā}.) Narrow? (\textbf{\ur{تنگ}}, \emph{tang}.) Deep? (\textbf{\ur{گہرا}}, \emph{gahrā}.) Famous? (\textbf{\ur{مشہور}}, \emph{mashhūr}.) Brave? (\textbf{\ur{بہادر}}, \emph{bahādur}.) Clever? (\textbf{\ur{ہوشیار}}, \emph{hoshiyār}.) Tired? (\textbf{\ur{تھکا}}, \emph{thakā}.) Hungry? (\textbf{\ur{بھوکا}}, \emph{bhūkā}.) Thirsty? (\textbf{\ur{پیاسا}}, \emph{pyāsā}.) Fresh? (\textbf{\ur{تازہ}}, \emph{tāzā}.) Raw? (\textbf{\ur{کچا}}, \emph{kachchā}.) Ripe? (\textbf{\ur{پکا}}, \emph{pakkā}.) Pink? (\textbf{\ur{گلابی}}, \emph{gulābī}.) Purple? (\textbf{\ur{جامنی}}, \emph{jāmnī}.) Golden? (\textbf{\ur{سنہرا}}, \emph{sunahrā}.) Closed? (\textbf{\ur{بند}}, \emph{band}.) Open? (\textbf{\ur{کھلا}}, \emph{khulā}.) To wake up? (\textbf{\ur{جاگنا}}, \emph{jāgnā}.)
\end{tcolorbox}
\section[(dialogue)]{Verbs — a second pass}

\label{lesson:UR-R142-second-pass-verbs}



\begin{quote}
Say the words for to fall and to fly.
\end{quote}

\subsection*{Your turn}
Say these aloud:

\begin{itemize}
\raggedright
  \item to fall, to fly, to swim, to play, to touch
  \item to make, to break, to burn, to run away, to call
  \item to win, to lose, to teach, to explain, to ask for
  \item to stay, to come back, to arrive, to come out, to climb
  \item to get off, to fill, to hide, to be saved, to grow
\end{itemize}

\begin{tcolorbox}[breakable,colback=teal!4,colframe=teal!35!black,title={Before you move on}]
To fall? (\textbf{\ur{گرنا}}, \emph{girnā}.) To fly? (\textbf{\ur{اڑنا}}, \emph{uṛnā}.) To swim? (\textbf{\ur{تیرنا}}, \emph{tairnā}.) To play? (\textbf{\ur{بجانا}}, \emph{bajānā}.) To touch? (\textbf{\ur{چھونا}}, \emph{chhūnā}.) To make? (\textbf{\ur{بنانا}}, \emph{banānā}.) To break? (\textbf{\ur{توڑنا}}, \emph{toṛnā}.) To burn? (\textbf{\ur{جلانا}}, \emph{jalānā}.) To run away? (\textbf{\ur{بھاگنا}}, \emph{bhāgnā}.) To call? (\textbf{\ur{بلانا}}, \emph{bulānā}.) To win? (\textbf{\ur{جیتنا}}, \emph{jītnā}.) To lose? (\textbf{\ur{ہارنا}}, \emph{hārnā}.) To teach? (\textbf{\ur{سکھانا}}, \emph{sikhānā}.) To explain? (\textbf{\ur{سمجھانا}}, \emph{samjhānā}.) To ask for? (\textbf{\ur{مانگنا}}, \emph{māṅgnā}.) To stay? (\textbf{\ur{ٹھہرنا}}, \emph{ṭhaharnā}.) To come back? (\textbf{\ur{لوٹنا}}, \emph{lauṭnā}.) To arrive? (\textbf{\ur{پہنچنا}}, \emph{pahuṅchnā}.) To come out? (\textbf{\ur{نکلنا}}, \emph{nikalnā}.) To climb? (\textbf{\ur{چڑھنا}}, \emph{chaṛhnā}.) To get off? (\textbf{\ur{اترنا}}, \emph{utarnā}.) To fill? (\textbf{\ur{بھرنا}}, \emph{bharnā}.) To hide? (\textbf{\ur{چھپنا}}, \emph{chhupnā}.) To be saved? (\textbf{\ur{بچنا}}, \emph{bachnā}.) To grow? (\textbf{\ur{بڑھنا}}, \emph{baṛhnā}.)
\end{tcolorbox}
\section[(timed A1 writing)]{Your first complete A1 writing paper}

\label{lesson:UR-W10-a1-timed-production}



\begin{quote}
This is an attempt, not a model. You will write the answers, and this page will not show you any copyable answer.

Put away Roman Urdu, Devanagari, dictionaries, translators, and spell-checkers. Set one timer for \textbf{20 minutes}. Keep it running through both tasks. Write in Urdu script; either approved Nastaliq or accessibility Naskh is fine. Ordinary legible handwriting is enough: imitation Nastaliq calligraphy is not scored.
\end{quote}

\subsection*{Writing — timed assessment production}
\#\#\# Task 1 — practical form: 7 minutes, 40 points

Complete \textbf{all six fields in Urdu script} for a new neighbourhood activity group. Short answers are enough.

\begin{enumerate}
\raggedright
  \item Your name
  \item Your city or village
  \item A language you speak
  \item One place you like
  \item One drink or food you like
  \item One everyday action or kind of work you do
\end{enumerate}

At seven minutes, move on even if a field is unfinished. Do not restart the timer and do not polish the form with time meant for the message.

\#\#\# Task 2 — message: 13 minutes, 60 points

\textbf{Reader:} Amina. \textbf{Purpose:} help Amina recognise you and arrange your first meeting tomorrow.

Write \textbf{30–40 words in Urdu script}. Greet Amina, introduce yourself, say where you live and which language you speak, mention one place and one drink or food you like, ask where Amina lives, and close politely. The instructions give content, not sentences: choose and connect the Urdu yourself.

When the 20-minute timer rings, \textbf{stop}, even if a line is unfinished. Record the number of completed form fields and the message word count. An honest timed attempt is more useful than a perfect untimed one.

\begin{tcolorbox}[breakable,colback=teal!4,colframe=teal!35!black,title={Before you move on}]
Only after you have stopped, review the same attempt. Do not rewrite it yet. Mark one strength and one next step under each assessment criterion:

\begin{enumerate}
\raggedright
  \item \textbf{Task fulfilment and relevant content} — six filled fields; the named reader, meeting purpose, and requested message points are present.
  \item \textbf{Comprehensibility and organisation} — Amina can follow the message and see how its ideas connect.
  \item \textbf{Urdu vocabulary, grammar, and register} — familiar words and sentence patterns carry the intended meaning in a suitable everyday tone.
  \item \textbf{Urdu orthographic control} — Urdu letters, joining, spacing, and punctuation remain readable. Optional short vowels or diacritics earn no bonus and receive no penalty unless a tested contrast depends on them.
\end{enumerate}

Keep the paper. On another day, repeat the same two-task shape with new personal details. The goal is not to memorise this prompt; it is to make a complete Urdu attempt calmly inside the real limit.
\end{tcolorbox}
